\section{Classical Mergeable-Sketch Parity: HyperLogLog}\label{sec:hll-parity}

Section~\ref{sec:markov-walkthrough} showed that local supervision can
move down the tree on a task whose oracle defeats unordered merges.
We now ask the converse: when the oracle \emph{is} a mergeable
sketch's query, does C-TreePO recover the classical sketch exactly? The answer is yes for
HyperLogLog \citep{FlajoletEtAl2007,HeuleEtAl2013}, the canonical
sketch for distinct counting. Running native HLL through the unified
\texttt{fit()} pipeline at $n_{\mathrm{leaves}} \in \{1, 2, 4, 8, 16\}$
produces register vectors that are byte-identical to the flat HLL on
every cell of the precision-by-leaf-count grid, and the Apache
DataSketches reference agrees with the flat run within twice the
relative standard error. This is the empirical face of
Proposition~\ref{prop:mergeable-reduction}'s state-level collapse:
when $g$ serializes the HLL register state and the merge is
register-wise max, the tree reduction inherits the sketch's
algebraic guarantees verbatim, with no learning required. The Hamlet
analogue is where the same construction earns its place outside
synthetic data: each scene emits the set of speaker pairs sharing
the stage; each act merges by set-union; HLL's register array
approximates union cardinality without storing the union set.

The same pipeline can swap its oracle from analytic distinct count
to the classical HLL's own scoring head, then ask whether a
\emph{learned} $g$ can recover the same behavior under local-law
supervision. Table~\ref{tab:classical-parity-hll} summarizes the four methods:
classical-native is byte-exact; classical-DataSketches is estimate-
exact within RSE; the end-to-end learned $g{+}f$ tracks the classical
sketch within a small multiple of the RSE floor and \emph{improves}
as the tree gets deeper; and a constrained learned-$g$-with-classical-
$f$ variant degrades with $L$, quantifying the cost of forcing a
learned operator to live inside the classical register-space readout.
The pattern is consistent: the projection onto the local-law subspace
is where C-TreePO matches the classical sketch, and the learned
representations gain a degree of freedom the byte-exact route does
not. Appendix~\ref{app:classical-parity} carries the protocol, the
sweep grid, the per-cell numbers, the full precision-by-leaf-count
parity figure, the register-space constraint analysis, the broader
\texttt{treepo-bench classical-sketches} discussion, and the
Count-Min/Frequent-Items taxonomy sidebar (these frequency sketches
are associative and commutative but not idempotent, so the C2
machinery applies differently than for HLL).

\begin{table}[t]
    \centering
    \small
    % Auto-generated by unified_g_v1.sketch.classical_parity_report; do not edit.
\begin{tabular}{llrrrrrrrrrrrrr}
\toprule
method & oracle & $p$ & $L$ & RSE & $\overline{|\Delta_{rel}|}$ & 95\% CI & $\max|\Delta|$ & state$=$ & bytes$=$ & rel MAE & C1 & C3 & wall$\times$ & $n$ \\
\midrule
classical\_datasketches & analytic & 11 & 1 & 0.02298 & 0 & 0 & 0 & 1 & 1 & 0.01133 & 15.14 & 0 & 0.8513 & 3 \\
classical\_datasketches & analytic & 11 & 16 & 0.02298 & 0.008052 & 0.0008999 & 13.05 & 1 & 0 & 0.01322 & 0.07798 & 8.668 & 1.832 & 3 \\
classical\_datasketches & analytic & 11 & 2 & 0.02298 & 0.008025 & 0.0008882 & 13.05 & 1 & 0 & 0.01319 & 7.093 & 17.24 & 1.134 & 3 \\
classical\_datasketches & analytic & 11 & 4 & 0.02298 & 0.008054 & 0.000927 & 13.07 & 1 & 0 & 0.01322 & 2.854 & 12.88 & 1.319 & 3 \\
classical\_datasketches & analytic & 11 & 8 & 0.02298 & 0.008053 & 0.0009013 & 13.05 & 1 & 0 & 0.01317 & 0.885 & 10.13 & 1.577 & 3 \\
classical\_datasketches & analytic & 13 & 1 & 0.01149 & 0 & 0 & 0 & 1 & 1 & 0.002884 & 5.227 & 0 & 1.032 & 3 \\
classical\_datasketches & analytic & 13 & 16 & 0.01149 & 0.002351 & 0.0008971 & 5.651 & 1 & 0 & 0.003686 & 3.394e-05 & 1.994 & 1.572 & 3 \\
classical\_datasketches & analytic & 13 & 2 & 0.01149 & 0.002365 & 0.0008783 & 5.53 & 1 & 0 & 0.003722 & 1.378 & 7.073 & 1.221 & 3 \\
classical\_datasketches & analytic & 13 & 4 & 0.01149 & 0.002415 & 0.0007943 & 5.538 & 1 & 0 & 0.003728 & 0.1452 & 3.593 & 1.376 & 3 \\
classical\_datasketches & analytic & 13 & 8 & 0.01149 & 0.002327 & 0.000767 & 5.452 & 1 & 0 & 0.003675 & 0.0009903 & 2.417 & 1.531 & 3 \\
classical\_datasketches & analytic & 7 & 1 & 0.09192 & 0 & 0 & 0 & 1 & 1 & 0.05798 & 73.93 & 0 & 1.053 & 3 \\
classical\_datasketches & analytic & 7 & 16 & 0.09192 & 0.04028 & 0.004683 & 59.86 & 1 & 0 & 0.07586 & 4.539 & 54.61 & 1.321 & 3 \\
classical\_datasketches & analytic & 7 & 2 & 0.09192 & 0.04028 & 0.004683 & 59.86 & 1 & 0 & 0.07586 & 35.33 & 102.8 & 1.073 & 3 \\
classical\_datasketches & analytic & 7 & 4 & 0.09192 & 0.04028 & 0.004683 & 59.86 & 1 & 0 & 0.07586 & 18.97 & 73.53 & 1.113 & 3 \\
classical\_datasketches & analytic & 7 & 8 & 0.09192 & 0.04028 & 0.004683 & 59.86 & 1 & 0 & 0.07586 & 9.047 & 61.06 & 0.9858 & 3 \\
classical\_datasketches & analytic & 9 & 1 & 0.04596 & 0 & 0 & 0 & 1 & 1 & 0.02289 & 29 & 0 & 1.115 & 3 \\
classical\_datasketches & analytic & 9 & 16 & 0.04596 & 0.01965 & 0.001969 & 31.12 & 1 & 0 & 0.02862 & 1.82 & 22.09 & 1.825 & 3 \\
classical\_datasketches & analytic & 9 & 2 & 0.04596 & 0.01965 & 0.001969 & 31.12 & 1 & 0 & 0.02862 & 15.27 & 39.39 & 1.085 & 3 \\
classical\_datasketches & analytic & 9 & 4 & 0.04596 & 0.01965 & 0.001969 & 31.12 & 1 & 0 & 0.02862 & 7.661 & 29.39 & 1.177 & 3 \\
classical\_datasketches & analytic & 9 & 8 & 0.04596 & 0.01965 & 0.001969 & 31.12 & 1 & 0 & 0.02862 & 3.956 & 24.61 & 1.321 & 3 \\
classical\_datasketches & hll\_reference & 11 & 1 & 0.02298 & 0 & 0 & 0 & 1 & 1 & 0 & 0 & 0 & 1.041 & 3 \\
classical\_datasketches & hll\_reference & 11 & 16 & 0.02298 & 0.008052 & 0.0008999 & 13.05 & 1 & 0 & 0.008052 & 0 & 2.508 & 1.83 & 3 \\
classical\_datasketches & hll\_reference & 11 & 2 & 0.02298 & 0.008025 & 0.0008882 & 13.05 & 1 & 0 & 0.008025 & 0 & 10.68 & 1.141 & 3 \\
classical\_datasketches & hll\_reference & 11 & 4 & 0.02298 & 0.008054 & 0.000927 & 13.07 & 1 & 0 & 0.008054 & 0 & 7.789 & 1.322 & 3 \\
classical\_datasketches & hll\_reference & 11 & 8 & 0.02298 & 0.008053 & 0.0009013 & 13.05 & 1 & 0 & 0.008053 & 0 & 5.765 & 1.513 & 3 \\
classical\_datasketches & hll\_reference & 13 & 1 & 0.01149 & 0 & 0 & 0 & 1 & 1 & 0 & 0 & 0 & 1.047 & 3 \\
classical\_datasketches & hll\_reference & 13 & 16 & 0.01149 & 0.002351 & 0.0008971 & 5.651 & 1 & 0 & 0.002351 & 0 & 0.6059 & 1.678 & 3 \\
classical\_datasketches & hll\_reference & 13 & 2 & 0.01149 & 0.002365 & 0.0008783 & 5.53 & 1 & 0 & 0.002365 & 0 & 4.251 & 1.186 & 3 \\
classical\_datasketches & hll\_reference & 13 & 4 & 0.01149 & 0.002415 & 0.0007943 & 5.538 & 1 & 0 & 0.002415 & 0 & 1.858 & 1.379 & 3 \\
classical\_datasketches & hll\_reference & 13 & 8 & 0.01149 & 0.002327 & 0.000767 & 5.452 & 1 & 0 & 0.002327 & 0 & 0.9663 & 1.535 & 3 \\
classical\_datasketches & hll\_reference & 7 & 1 & 0.09192 & 0 & 0 & 0 & 1 & 1 & 0 & 0 & 0 & 1.054 & 3 \\
classical\_datasketches & hll\_reference & 7 & 16 & 0.09192 & 0.04028 & 0.004683 & 59.86 & 1 & 0 & 0.04028 & 0 & 32.12 & 1.308 & 3 \\
classical\_datasketches & hll\_reference & 7 & 2 & 0.09192 & 0.04028 & 0.004683 & 59.86 & 1 & 0 & 0.04028 & 0 & 53.52 & 1.078 & 3 \\
classical\_datasketches & hll\_reference & 7 & 4 & 0.09192 & 0.04028 & 0.004683 & 59.86 & 1 & 0 & 0.04028 & 0 & 41.47 & 1.121 & 3 \\
classical\_datasketches & hll\_reference & 7 & 8 & 0.09192 & 0.04028 & 0.004683 & 59.86 & 1 & 0 & 0.04028 & 0 & 35.49 & 1.102 & 3 \\
classical\_datasketches & hll\_reference & 9 & 1 & 0.04596 & 0 & 0 & 0 & 1 & 1 & 0 & 0 & 0 & 1.052 & 3 \\
classical\_datasketches & hll\_reference & 9 & 16 & 0.04596 & 0.01965 & 0.001969 & 31.12 & 1 & 0 & 0.01965 & 0 & 14.5 & 1.568 & 3 \\
classical\_datasketches & hll\_reference & 9 & 2 & 0.04596 & 0.01965 & 0.001969 & 31.12 & 1 & 0 & 0.01965 & 0 & 27.06 & 0.8283 & 3 \\
classical\_datasketches & hll\_reference & 9 & 4 & 0.04596 & 0.01965 & 0.001969 & 31.12 & 1 & 0 & 0.01965 & 0 & 19.52 & 1.166 & 3 \\
classical\_datasketches & hll\_reference & 9 & 8 & 0.04596 & 0.01965 & 0.001969 & 31.12 & 1 & 0 & 0.01965 & 0 & 16.1 & 1.314 & 3 \\
classical\_native & analytic & 11 & 1 & 0.02298 & 0 & 0 & 0 & 1 & 1 & 0.01461 & 19.47 & 0 & 1.018 & 3 \\
classical\_native & analytic & 11 & 16 & 0.02298 & 0 & 0 & 0 & 1 & 1 & 0.01461 & 1.441 & 10.92 & 1.135 & 3 \\
classical\_native & analytic & 11 & 2 & 0.02298 & 0 & 0 & 0 & 1 & 1 & 0.01461 & 10.44 & 19.47 & 0.9988 & 3 \\
classical\_native & analytic & 11 & 4 & 0.02298 & 0 & 0 & 0 & 1 & 1 & 0.01461 & 5.081 & 14.54 & 1.041 & 3 \\
classical\_native & analytic & 11 & 8 & 0.02298 & 0 & 0 & 0 & 1 & 1 & 0.01461 & 2.753 & 12.23 & 1.061 & 3 \\
classical\_native & analytic & 13 & 1 & 0.01149 & 0 & 0 & 0 & 1 & 1 & 0.008445 & 9.396 & 0 & 0.9912 & 3 \\
classical\_native & analytic & 13 & 16 & 0.01149 & 0 & 0 & 0 & 1 & 1 & 0.008445 & 0.7089 & 5.498 & 1.094 & 3 \\
classical\_native & analytic & 13 & 2 & 0.01149 & 0 & 0 & 0 & 1 & 1 & 0.008445 & 5.149 & 9.396 & 1.025 & 3 \\
classical\_native & analytic & 13 & 4 & 0.01149 & 0 & 0 & 0 & 1 & 1 & 0.008445 & 2.723 & 7.212 & 1.017 & 3 \\
classical\_native & analytic & 13 & 8 & 0.01149 & 0 & 0 & 0 & 1 & 1 & 0.008445 & 1.405 & 6.127 & 1.074 & 3 \\
classical\_native & analytic & 7 & 1 & 0.09192 & 0 & 0 & 0 & 1 & 1 & 0.06938 & 84.16 & 0 & 0.9672 & 3 \\
classical\_native & analytic & 7 & 16 & 0.09192 & 0 & 0 & 0 & 1 & 1 & 0.06938 & 6.93 & 51 & 1.139 & 3 \\
classical\_native & analytic & 7 & 2 & 0.09192 & 0 & 0 & 0 & 1 & 1 & 0.06938 & 49.8 & 84.16 & 1.02 & 3 \\
classical\_native & analytic & 7 & 4 & 0.09192 & 0 & 0 & 0 & 1 & 1 & 0.06938 & 24.07 & 65.44 & 1.054 & 3 \\
classical\_native & analytic & 7 & 8 & 0.09192 & 0 & 0 & 0 & 1 & 1 & 0.06938 & 13.36 & 56.09 & 1.094 & 3 \\
classical\_native & analytic & 9 & 1 & 0.04596 & 0 & 0 & 0 & 1 & 1 & 0.03318 & 41.96 & 0 & 0.9943 & 3 \\
classical\_native & analytic & 9 & 16 & 0.04596 & 0 & 0 & 0 & 1 & 1 & 0.03318 & 2.744 & 24.37 & 1.12 & 3 \\
classical\_native & analytic & 9 & 2 & 0.04596 & 0 & 0 & 0 & 1 & 1 & 0.03318 & 23.93 & 41.96 & 1.05 & 3 \\
classical\_native & analytic & 9 & 4 & 0.04596 & 0 & 0 & 0 & 1 & 1 & 0.03318 & 11.5 & 31.98 & 1.07 & 3 \\
classical\_native & analytic & 9 & 8 & 0.04596 & 0 & 0 & 0 & 1 & 1 & 0.03318 & 5.598 & 27.01 & 1.069 & 3 \\
classical\_native & hll\_reference & 11 & 1 & 0.02298 & 0 & 0 & 0 & 1 & 1 & 0 & 0 & 0 & 1.009 & 3 \\
classical\_native & hll\_reference & 11 & 16 & 0.02298 & 0 & 0 & 0 & 1 & 1 & 0 & 0 & 0 & 1.096 & 3 \\
classical\_native & hll\_reference & 11 & 2 & 0.02298 & 0 & 0 & 0 & 1 & 1 & 0 & 0 & 0 & 1.017 & 3 \\
classical\_native & hll\_reference & 11 & 4 & 0.02298 & 0 & 0 & 0 & 1 & 1 & 0 & 0 & 0 & 1.066 & 3 \\
classical\_native & hll\_reference & 11 & 8 & 0.02298 & 0 & 0 & 0 & 1 & 1 & 0 & 0 & 0 & 1.03 & 3 \\
classical\_native & hll\_reference & 13 & 1 & 0.01149 & 0 & 0 & 0 & 1 & 1 & 0 & 0 & 0 & 0.9949 & 3 \\
classical\_native & hll\_reference & 13 & 16 & 0.01149 & 0 & 0 & 0 & 1 & 1 & 0 & 0 & 0 & 1.08 & 3 \\
classical\_native & hll\_reference & 13 & 2 & 0.01149 & 0 & 0 & 0 & 1 & 1 & 0 & 0 & 0 & 1.007 & 3 \\
classical\_native & hll\_reference & 13 & 4 & 0.01149 & 0 & 0 & 0 & 1 & 1 & 0 & 0 & 0 & 1.033 & 3 \\
classical\_native & hll\_reference & 13 & 8 & 0.01149 & 0 & 0 & 0 & 1 & 1 & 0 & 0 & 0 & 1.011 & 3 \\
classical\_native & hll\_reference & 7 & 1 & 0.09192 & 0 & 0 & 0 & 1 & 1 & 0 & 0 & 0 & 0.9832 & 3 \\
classical\_native & hll\_reference & 7 & 16 & 0.09192 & 0 & 0 & 0 & 1 & 1 & 0 & 0 & 0 & 1.158 & 3 \\
classical\_native & hll\_reference & 7 & 2 & 0.09192 & 0 & 0 & 0 & 1 & 1 & 0 & 0 & 0 & 1.05 & 3 \\
classical\_native & hll\_reference & 7 & 4 & 0.09192 & 0 & 0 & 0 & 1 & 1 & 0 & 0 & 0 & 1.069 & 3 \\
classical\_native & hll\_reference & 7 & 8 & 0.09192 & 0 & 0 & 0 & 1 & 1 & 0 & 0 & 0 & 1.11 & 3 \\
classical\_native & hll\_reference & 9 & 1 & 0.04596 & 0 & 0 & 0 & 1 & 1 & 0 & 0 & 0 & 0.9869 & 3 \\
classical\_native & hll\_reference & 9 & 16 & 0.04596 & 0 & 0 & 0 & 1 & 1 & 0 & 0 & 0 & 1.146 & 3 \\
classical\_native & hll\_reference & 9 & 2 & 0.04596 & 0 & 0 & 0 & 1 & 1 & 0 & 0 & 0 & 1.052 & 3 \\
classical\_native & hll\_reference & 9 & 4 & 0.04596 & 0 & 0 & 0 & 1 & 1 & 0 & 0 & 0 & 1.032 & 3 \\
classical\_native & hll\_reference & 9 & 8 & 0.04596 & 0 & 0 & 0 & 1 & 1 & 0 & 0 & 0 & 1.076 & 3 \\
learned\_g & hll\_reference & 11 & 1 & 0.02298 & -- & -- & -- & -- & -- & 0.3265 & 66.2 & 0 & -- & 3 \\
learned\_g & hll\_reference & 11 & 16 & 0.02298 & -- & -- & -- & -- & -- & 0.8796 & 3.599 & 95.57 & -- & 3 \\
learned\_g & hll\_reference & 11 & 2 & 0.02298 & -- & -- & -- & -- & -- & 0.5141 & 31.86 & 105.6 & -- & 3 \\
learned\_g & hll\_reference & 11 & 4 & 0.02298 & -- & -- & -- & -- & -- & 0.4717 & 13.91 & 70.56 & -- & 3 \\
learned\_g & hll\_reference & 11 & 8 & 0.02298 & -- & -- & -- & -- & -- & 0.7468 & 6.413 & 87.07 & -- & 3 \\
learned\_g & hll\_reference & 13 & 1 & 0.01149 & -- & -- & -- & -- & -- & 0.2866 & 58.25 & 0 & -- & 3 \\
learned\_g & hll\_reference & 13 & 16 & 0.01149 & -- & -- & -- & -- & -- & 0.8893 & 3.847 & 98.03 & -- & 3 \\
learned\_g & hll\_reference & 13 & 2 & 0.01149 & -- & -- & -- & -- & -- & 0.4998 & 35.13 & 103.7 & -- & 3 \\
learned\_g & hll\_reference & 13 & 4 & 0.01149 & -- & -- & -- & -- & -- & 0.7436 & 16.13 & 107.5 & -- & 3 \\
learned\_g & hll\_reference & 13 & 8 & 0.01149 & -- & -- & -- & -- & -- & 0.7201 & 7.373 & 82.79 & -- & 3 \\
learned\_g & hll\_reference & 7 & 1 & 0.09192 & -- & -- & -- & -- & -- & 0.3292 & 65.61 & 0 & -- & 3 \\
learned\_g & hll\_reference & 7 & 16 & 0.09192 & -- & -- & -- & -- & -- & 0.186 & 8.146 & 27.95 & -- & 3 \\
learned\_g & hll\_reference & 7 & 2 & 0.09192 & -- & -- & -- & -- & -- & 0.2623 & 29.34 & 52.37 & -- & 3 \\
learned\_g & hll\_reference & 7 & 4 & 0.09192 & -- & -- & -- & -- & -- & 0.1863 & 14.18 & 33.83 & -- & 3 \\
learned\_g & hll\_reference & 7 & 8 & 0.09192 & -- & -- & -- & -- & -- & 0.1778 & 8.486 & 28.07 & -- & 3 \\
learned\_g & hll\_reference & 9 & 1 & 0.04596 & -- & -- & -- & -- & -- & 0.2654 & 53.36 & 0 & -- & 3 \\
learned\_g & hll\_reference & 9 & 16 & 0.04596 & -- & -- & -- & -- & -- & 0.4396 & 4.974 & 53.75 & -- & 3 \\
learned\_g & hll\_reference & 9 & 2 & 0.04596 & -- & -- & -- & -- & -- & 0.2657 & 30.51 & 53.77 & -- & 3 \\
learned\_g & hll\_reference & 9 & 4 & 0.04596 & -- & -- & -- & -- & -- & 0.2919 & 16.34 & 44.86 & -- & 3 \\
learned\_g & hll\_reference & 9 & 8 & 0.04596 & -- & -- & -- & -- & -- & 0.2835 & 8.893 & 39.25 & -- & 3 \\
learned\_fg & hll\_reference & 11 & 1 & 0.02298 & -- & -- & -- & -- & -- & 0.2963 & 60.48 & 0 & -- & 3 \\
learned\_fg & hll\_reference & 11 & 16 & 0.02298 & -- & -- & -- & -- & -- & 0.2481 & 5.789 & 31.28 & -- & 3 \\
learned\_fg & hll\_reference & 11 & 2 & 0.02298 & -- & -- & -- & -- & -- & 0.2488 & 29.12 & 50.53 & -- & 3 \\
learned\_fg & hll\_reference & 11 & 4 & 0.02298 & -- & -- & -- & -- & -- & 0.2465 & 14.89 & 39.7 & -- & 3 \\
learned\_fg & hll\_reference & 11 & 8 & 0.02298 & -- & -- & -- & -- & -- & 0.223 & 8.897 & 33.8 & -- & 3 \\
learned\_fg & hll\_reference & 13 & 1 & 0.01149 & -- & -- & -- & -- & -- & 0.2649 & 53.93 & 0 & -- & 3 \\
learned\_fg & hll\_reference & 13 & 16 & 0.01149 & -- & -- & -- & -- & -- & 0.2489 & 7.425 & 31.66 & -- & 3 \\
learned\_fg & hll\_reference & 13 & 2 & 0.01149 & -- & -- & -- & -- & -- & 0.2657 & 30.16 & 54.39 & -- & 3 \\
learned\_fg & hll\_reference & 13 & 4 & 0.01149 & -- & -- & -- & -- & -- & 0.2366 & 13.55 & 36.4 & -- & 3 \\
learned\_fg & hll\_reference & 13 & 8 & 0.01149 & -- & -- & -- & -- & -- & 0.1825 & 7.409 & 34.55 & -- & 3 \\
learned\_fg & hll\_reference & 7 & 1 & 0.09192 & -- & -- & -- & -- & -- & 0.298 & 59.83 & 0 & -- & 3 \\
learned\_fg & hll\_reference & 7 & 16 & 0.09192 & -- & -- & -- & -- & -- & 0.2491 & 6.522 & 34.85 & -- & 3 \\
learned\_fg & hll\_reference & 7 & 2 & 0.09192 & -- & -- & -- & -- & -- & 0.2726 & 29.94 & 53.97 & -- & 3 \\
learned\_fg & hll\_reference & 7 & 4 & 0.09192 & -- & -- & -- & -- & -- & 0.2514 & 13.81 & 38 & -- & 3 \\
learned\_fg & hll\_reference & 7 & 8 & 0.09192 & -- & -- & -- & -- & -- & 0.2326 & 7.277 & 37.82 & -- & 3 \\
learned\_fg & hll\_reference & 9 & 1 & 0.04596 & -- & -- & -- & -- & -- & 0.2949 & 59.31 & 0 & -- & 3 \\
learned\_fg & hll\_reference & 9 & 16 & 0.04596 & -- & -- & -- & -- & -- & 0.2485 & 6.742 & 36.4 & -- & 3 \\
learned\_fg & hll\_reference & 9 & 2 & 0.04596 & -- & -- & -- & -- & -- & 0.2687 & 30.47 & 54.14 & -- & 3 \\
learned\_fg & hll\_reference & 9 & 4 & 0.04596 & -- & -- & -- & -- & -- & 0.2599 & 14.64 & 40.82 & -- & 3 \\
learned\_fg & hll\_reference & 9 & 8 & 0.04596 & -- & -- & -- & -- & -- & 0.2054 & 10.76 & 41.49 & -- & 3 \\
\bottomrule
\end{tabular}

    \caption{HLL parity summary across the \texttt{(method, oracle, precision, $L$)} grid. \texttt{RSE} is the HLL theoretical relative standard error $1.04/\sqrt{2^p}$. \texttt{state=} and \texttt{bytes=} are per-cell equality rates for the classical rows; \texttt{rel MAE} is each method's relative mean absolute error against the oracle; \texttt{C1} and \texttt{C3} are the per-leaf and per-merge MAE; \texttt{wall$\times$} is the tree-to-flat wall-time ratio. Full table with all seeds and oracle-kind variants in Appendix~\ref{app:classical-parity}.}
    \label{tab:classical-parity-hll}
\end{table}
